\section{\sys Design}
\label{sec:design}

\sys is a recovery layer beneath an unmodified GUI agent. It intercepts the agent's actions, records effects before they commit, and after any interruption
decides each ambiguous effect with the rules of \S\ref{sec:theory}. Its parts are an effect classifier, the ledger, the nonce writer, the probe families with their
calibration store, and the \alg planner with its saga layer (Figure~\ref{fig:arch} in Appendix~\ref{app:extra}).

\subsection{Where it runs, and the ledger}
\label{sec:design-ledger}
The ledger must outlive the agent. In the host topology (our testbed), the controller is a host process separate from the agent, which drives the
emulator over adb, so killing the agent cannot erase the ledger. In the on-device topology, the controller is an Android foreground service in its own process, with the ledger
in app-private storage. The low-memory killer can still reap it, but safety does not depend on it surviving: the write-ahead invariant below holds across its death, and it resolves
outstanding records when it restarts. We have built the host topology's decision core and ledger, and the on-device service is design only.

Agents with structured plans label commit points directly~\cite{tian2025agentprog}. Otherwise a conservative classifier matches the target's text, content
description or resource id against a per-locale lexicon (send, pay, save, post, book, confirm, delete), or a VLM flags it~\cite{lee2025verisafe,sun2025ossentinel}. A false
positive costs one ledger write, and a false negative falls back to the agent's native recovery.

In the write-ahead ledger, each effect has a record \{\textit{id}, workflow, step, app package and version, description, nonce, history digest, attempt, state, $t_0$\}. States are
\textsf{PLANNED} $\to$ \textsf{ISSUED} $\to$ \{\textsf{COMMITTED}, \textsf{ABSENT}, \textsf{PENDING}, \textsf{ASKED}\}, with \textsf{ABSENT} $\to$ \textsf{ISSUED} on retry
(attempt$+1$, same nonce) and \textsf{COMMITTED} $\to$ \textsf{COMPENSATING} $\to$ \textsf{COMPENSATED} when a duplicate is detected. The invariant is write-ahead:
\textsf{ISSUED} is committed in SQLite (WAL, \texttt{synchronous=FULL}) before the commit tap, and on restart every \textsf{ISSUED} record is resolved before the agent continues.


\subsection{Content-carried nonce}
\label{sec:design-nonce}
A per-app field map names a free-text field that the app stores and displays: a private note on an event or contact, a memo on a transfer, a trailing tag on a message.
Before the commit tap, \sys types \texttt{gx:} and six Crockford base-32 symbols (30 bits; the alphabet omits I, L, O, U) into that field, and a retry reuses the tag.
The nonce probe reads the app's sent/history view through the accessibility tree and through in-app search, matches outgoing items of the verified account within edit distance~1,
and ignores drafts, failed entries and notifications (Lemma~\ref{lem:nonce}). With $M{=}40$ visible tokens this gives $\epsilon_0\le\NonceEpsZero$ per view (Table~\ref{tab:nonce}).
By default \sys uses only fields that the counterparty cannot see. Counterparty-visible fields (an SMS or chat body, a payment memo printed on a statement) need per-app opt-in, change
the content of the effect, and break graders that compare content exactly, such as AndroidWorld's SMS and Markor checks, unless they strip \texttt{gx:} tags.

\subsection{Probe families and calibration}
\label{sec:design-probes}
Each probe is a short navigation script plus a reader, grouped into four families.
(1)~Hop-1 counts outbox and sent-list entries created after $t_0$ that match the ledger description or nonce, in two views (outbox list, conversation status).
(2)~Nonce has two views as above.
(3)~Digest compares against a digest of matching history items that the ledger stores before the tap. After recovery the probe counts matching items in the history list and in
search results and reports whether a new one appeared. This answers look-alikes without a nonce, since last month's transfer is already in the digest.
(4)~VLM asks a yes/no question about the ledger description, once each of the list, detail and search views. The VLM family is skipped when the digest shows a look-alike.
Hop~1 reads the app's Outbox screen through the accessibility tree, like every other probe, and never the app's database. An app whose status is drawn only with the
rendered history (Element) has no hop-1 family. The other three families read the history as the app renders it, and in-app search often queries the server. They are therefore one
evidence source, and a stale history after a restart can fool all three at once. A server-backed search view inherits the hop-2 lag, so it counts toward ``absent'' only after the wait.

\para{Calibration labels.} For each (app version, family, verdict level or region) the store holds one-sided Clopper-Pearson bounds (Lemma~\ref{lem:calib}(ii)) from calibration
effects, not reads. Lemma~\ref{lem:calib} needs effects drawn from the interrupted population, so labels come from fault-injected runs on test accounts or a staging backend.
The controller injects the fault classes of \S\ref{sec:impl} at the commit point, records the true outcome from the test backend, runs every family once per view and logs the lag $T_c$ from $r_e$. Bounds and lag tables are keyed by the observable fault class (Lemma~\ref{lem:calib}(ii)), so the lab's fault mix need not
match deployment's. An effect whose class the controller cannot tell uses the per-class maximum. Passive labels (confirmed sends, abandoned forms) are cheaper but miss cached history after a
kill, optimistic entries and restart lags (\S\ref{sec:eval-pop},~\S\ref{sec:eval-real}). A label costs one fault-injected trial on a test device. Our scripted label trials on Etar and DAVx5 (\S\ref{sec:eval-real}) took a median of \CertLabelMedOneS\,s for a committed
effect and \CertLabelMedZeroS\,s for an absent one. At that rate, \CertLabelsFive labels per class and outcome
for \CertClassesN classes take about \CertDeviceHoursFive device-hours, and \CertLagBootN{} lag entries per class about \CertDeviceHoursLagFive{} more.

\para{How many labels.} Validity holds at any sample size: small samples only widen the bounds, so the planner asks more often, not retries more often (Table~\ref{tab:nsize}).
Product certification needs a few hundred effects per stratum at $\alpha{=}1\%$, because per-family bounds multiply. Joint certification and each wait need about $\ln(1/\delta)/\alpha_{\mathrm{test}}$ samples even with zero events, which is thousands at $\alpha{=}1\%$.

Four different counts make up the label budget of our measured-parameter simulations (Table~\ref{tab:cert-budget}). The Element profile rests on \CertElemTrialsN emulator trials.
Each simulation draws \CertLabelsBase or \CertLabelsFive labels per fault class and outcome from that profile. Joint certification splits each set into a pilot that fixes the score
and a holdout that tests it, \CertPilotFive and \CertHoldoutFive labels at the larger size. Unless a row ties them to the label count, the lag tables hold \CertLagBootN entries per class, drawn by bootstrap from the \RsimLagSrcLo{} to \RsimLagSrcHi{} measured lags of each class, so none exceeds the largest measured lag.

The lag table is the costly part (Table~\ref{tab:cert-budget}). With as many lag entries as labels, the simulated joint mode asks on \CertTiedFiveJointAsk\% at \CertLabelsFive labels
instead of \RsimFiveJointAsk\%, and at \CertLabelsBase no wait certifies (larger budgets: Table~\ref{tab:cert-lagapp}). A lag tail that no label sees breaks the budget. When
\RsimTail\% of committed effects become visible only after every logged lag, product retries \CertTailFiveProdPv\% at \CertLabelsFive labels (CP \CertTailFiveProdPvHi\%). Joint keeps the
budget only at \CertLabelsBase labels, by asking on \RsimTailJointAsk\% of workflows, and tying the lag table to the labels does not help.


A single user produces nowhere near these counts per app version, so the store is pooled. Lag tables are pooled per app and backend across versions: lag
comes mostly from the network path and the app's sync, but the reader's navigation (most of the force-stop lag) depends on the UI and needs a fresh sample. Verdict bounds are pooled per app across versions and devices, with a shift test on a small sample of the new version.
When the test detects an increase, the planner falls back to that sample's own bounds, which means more asks, not fewer guarantees. The shift test is per family and level.
Pooling joint regions across versions needs a region-level test, which we have not built. The nonce's false-match bound is analytic. Pooling across devices assumes the devices
are exchangeable within a stratum (app version, locale, VLM), and per-device drift is not detected yet.

\subsection{The \alg planner}
\label{sec:design-planner}

\begin{algorithm}[t]
\caption{\alg: resolve one ghost effect $e$}
\label{alg:recover}
\small
\begin{algorithmic}[1]
\Require record $e$ (\textsf{ISSUED} at $t_0$, anchor $r_e$ of its fault class); budget $(\alpha_e,\beta_e)$ from the saga layer; calibrated bounds $q$; lag tables; idem: the app retries idempotently
\State $\alpha_t,\alpha_l,\beta_t\gets 0.5\alpha_e,\,0.4\alpha_e,\,0.9\beta_e$ \Comment{$\alpha_{\mathrm{cal}}$ via $\delta$}
\State $\tau_a,\tau_c\gets$ largest thresholds certified for $\alpha_t,\beta_t$ \Comment{Prop.~\ref{prop:family} or \ref{prop:joint}}
\State \Call{VerifyForeground}{$e$.app}; $P\gets1$; $d\gets\bot$; $\ell\gets\alpha_l$
\If{$e$.app has a hop-1 family}
  \State $\ell\gets\alpha_l/2$; $(s,r)\gets$ \Call{Readout}{hop-1} \Comment{entries after $t_0$}\label{alg:hop1}
  \If{$s=$ com $\wedge$ (idem $\vee$ local)} $P\gets q^{\mathrm{com}}_{\text{hop1}}(r)$; $d\gets$ com\label{alg:g2}
    \State $o\gets$ acked; \textbf{if} idem: $o\gets$ \Call{Monitor}{$e$} \Comment{bkgd., fallible}
    \State \textbf{if} $o=$ acked $\wedge\,P\le\tau_c$: \Return \Call{Mark}{$e$, \textsf{COMMITTED}}
    \State \textbf{if} $o\in\{$failed, stuck$\}$: $P\gets1$; $d\gets\bot$ \Comment{hand over, once}
  \ElsIf{$s=$ abs} $P\gets q^{\mathrm{abs}}_{\text{hop1}}(r)$; $d\gets$ abs; \textbf{if} $P\le\tau_a$ \Return \Call{Retry}{$e$}\label{alg:abs1}
  \EndIf
\EndIf
\If{$d\ne$ com \textbf{and} app not local}
  \State $w\gets$ \Call{CertifiedWait}{$e$.app, $e$.class, $d{=}$abs ? $\ell/q_h$ : $\ell$} \Comment{$q_h{=}\max_r q^{\mathrm{abs}}_{\text{hop1}}(r)$}\label{alg:cw}
  \State $D\gets r_e+w$ \Comment{class unknown: $\max_c(r_c+w_c)$}
  \State \textbf{if} an anchor $D$ needs is missing $\vee$ $D-\text{now}>W_{\max}$: \Return \Call{Ask}{$e$}\label{alg:wcap}
  \State \Call{WaitUntil}{$D$}\label{alg:wait}
\EndIf
\For{$f\in$ [nonce, digest, VLM] available for $e$}
  \State $(s,r)\gets$ \Call{Readout}{$f$} \Comment{one read per view}\label{alg:read}
  \State \textbf{if} $s=\bot$ \textbf{continue}; \textbf{if} $d\notin\{\bot,s\}$ \textbf{break}
  \State $P\gets P\cdot q^{s}_f(r)$; $d\gets s$
  \State \textbf{if} $s=$ abs $\wedge P\le\tau_a$: \Call{VerifyForeground}{}; \Return \Call{Retry}{$e$}
  \State \textbf{if} $s=$ com $\wedge P\le\tau_c$: \Return \Call{Mark}{$e$, \textsf{COMMITTED}}\label{alg:com}
\EndFor
\State \Return \Call{Ask}{$e$} \Comment{disagreement or not certifiable}\label{alg:ask}
\end{algorithmic}
\end{algorithm}

Algorithm~\ref{alg:recover} applies \S\ref{sec:theory} to one effect. \textsc{Readout} reads each view of a family once and returns the majority side and its count.
The thresholds come from enumerating stopping assignments (at most $4^4$ per app, precomputed) or, for joint certification, from the fixed-sequence test. \textsc{CertifiedWait} returns
the smallest logged lag $T_c$ of the effect's fault class whose CP upper bound on the exceedance probability is at most its argument. If the class is unknown, the
deadline $D$ is the latest over classes, and a missing anchor leads to an ask (\S\ref{sec:model}). A logged effect that never became visible counts as an exceedance at every wait. \textsc{Retry} reuses the nonce.

Two runs show the algorithm at work. When K-9 loses the SMTP reply to an applied send, its Outbox keeps a failed entry after $t_0$ (line~\ref{alg:hop1}). K-9 shows the same
entry when the server never applied the send, so line~\ref{alg:g2} trusts it only for apps with idempotent retries and local apps. No other family sees the mail, or a Sent
folder's lag exceeds $W_{\max}$ (line~\ref{alg:wcap}), so the planner asks (line~\ref{alg:ask}), as in the LLM-agent experiment (\S\ref{sec:eval-agent}). An interrupted
Element message has no hop-1 family. The planner waits until $r_e+w$ for its class (lines~\ref{alg:cw} and~\ref{alg:wait}) and reads the nonce views once (line~\ref{alg:read}).
It then marks the effect committed (line~\ref{alg:com}) or retries it with the same nonce once a product crosses its threshold, and asks otherwise. \S\ref{sec:eval-real} runs this path held out.

A hop-1 verdict starts the product on either side. The ``committed'' side of an app with idempotent retries, or of a local app, is therefore certified like any other declaration
(its bound enters $B_{\mathrm{com}}$), and a hop-1 ``absent'' verdict relaxes its own wait $w_1$, with half of the lag share (Eq.~\eqref{eq:lagsplit}). For an app with idempotent retries and an outbox entry,
\textsc{Monitor} marks the record \textsf{PENDING}, lets the agent continue and re-reads hop~1 in the background, and an acknowledged entry that hop~1 alone does not certify is corroborated end to end.
A failed entry, or one pending past a stuck limit (120\,s) or near $H_{\mathrm{dedup}}$, is handed over once to the end-to-end path, whose wait $w_0$ still runs from $r_e$ on the unconditional lag table, as Eq.~\eqref{eq:lagsplit} needs. Two entries mean a duplicate, which is compensated or asked about.

The wait is the price of the lag share (Theorem~\ref{thm:main}(a,b)). If the certified wait exceeds $W_{\max}$ (120\,s), or recovery was delayed while the device
dozed and no lag table exists for that state, the planner asks instead of waiting. Local-only apps have no hop~2, so their lag share is returned to the test share.

When the planner asks, the prompt names the effect, shows the pre-tap screenshot and the digest, and offers ``done'', ``not done'' and ``not sure''. An app that wins
the foreground can redirect a retry~\cite{qian2026rebinding}, so \textsc{VerifyForeground} checks package, activity and window before every retry, compensation or prompt.

A workflow's budget is split across effects by Eq.~\eqref{eq:kkt} in Appendix~\ref{app:proofs} with each step's strongest family as its probe; a duplicate of a reversible effect detected
later (two nonce matches, or two hop-1 entries) triggers the compensating effect as a new ledger record.
